\documentclass[11pt]{article}
\usepackage[margin=1in]{geometry}
\usepackage{amsmath,amssymb,booktabs,array,longtable}
\usepackage{graphicx}
\usepackage{hyperref}
\usepackage{enumitem}
\usepackage{titlesec}
\usepackage{listings}
\usepackage{xcolor}
\usepackage{float}   % <---- ADDED for [H] placement

\lstset{
    language=Python,
    basicstyle=\ttfamily\small,
    frame=single,
    numbers=left,
    breaklines=true,
    showstringspaces=false,
    keywordstyle=\color{blue},
    commentstyle=\color{green!60!black},
    stringstyle=\color{red}
}

\begin{document}

\begin{center}
\Large\textbf{Shiv Nadar University Chennai}\\[2mm]
\end{center}

\renewcommand{\arraystretch}{1.4}

\begin{tabular}{|p{3.5cm}|p{8cm}|p{2cm}|}
\hline
Degree \& Branch & B.Tech Artificial Intelligence \& Data Science & Semester V\\ \hline
Subject Code \& Name & CS3807 -- Deep Learning Laboratory & AY: 2026--27\\ \hline
\end{tabular}

\vspace{3mm}
\begin{center}
    {\Huge \textbf{Deep Learning Assignment Report }}
\end{center}

\begin{center}
\Large\textbf{Experiment 1}\\
\large\textbf{Implementation of a Single Layer Perceptron for Binary Classification}
\end{center}


%----------------------------------------------------------

\section*{1. Objective}

The objective of this experiment is to understand the concept of an artificial neuron and implement a Single Layer Perceptron from scratch for binary classification. Students will learn the perceptron learning algorithm, understand the importance of activation functions, visualize the learning process, and evaluate the classifier using a real-world dataset.

%----------------------------------------------------------

\section*{2. Tasks to be Performed}

\begin{enumerate}[leftmargin=*]
\item Download and understand the dataset.
\item Perform Exploratory Data Analysis (EDA).
\item Preprocess the dataset.
\item Implement a Single Layer Perceptron from scratch.
\item Train the model using the perceptron learning rule.
\item Evaluate the classifier using standard performance metrics.
\item Analyze the effect of different learning rates.
\item Compare the implementation with Scikit-learn's Perceptron (optional).
\end{enumerate}

%----------------------------------------------------------

\section*{3. Background Theory}

An Artificial Neuron is the fundamental building block of a neural network. It receives one or more input features, multiplies them by corresponding weights, adds a bias term, and applies an activation function to determine the output.

The Single Layer Perceptron is the earliest supervised learning algorithm capable of solving linearly separable binary classification problems.

\subsection*{Mathematical Formulation}

Weighted Sum

\[
z=\mathbf{w}^{T}\mathbf{x}+b
\]

where

\begin{itemize}
\item $\mathbf{x}$ : Input vector
\item $\mathbf{w}$ : Weight vector
\item $b$ : Bias
\item $z$ : Linear combination
\end{itemize}

Prediction

\[
\hat y=f(z)
\]

Weight Update Rule

\[
\mathbf{w}_{new}
=
\mathbf{w}_{old}
+
\eta
(y-\hat y)
\mathbf{x}
\]

Bias Update

\[
b_{new}
=
b_{old}
+
\eta
(y-\hat y)
\]

where $\eta$ denotes the learning rate.

%----------------------------------------------------------

\section*{4. Activation Functions}

An activation function determines the output of a neuron based on the weighted sum of its inputs. It introduces non-linearity, enabling neural networks to learn complex patterns. Although the perceptron uses only the Step Activation Function, modern deep learning models employ differentiable activation functions for efficient optimization through backpropagation.

\subsection*{Step Activation Function}

\[
f(z)=
\begin{cases}
1,&z\ge0\\
0,&z<0
\end{cases}
\]

The Step Activation Function produces a binary output and is therefore suitable only for linearly separable binary classification problems.

\begin{center}
\begin{tabular}{|p{2.5cm}|p{3cm}|p{6cm}|}
\hline
\textbf{Activation} & \textbf{Output Range} & \textbf{Typical Applications}\\
\hline
Step & $\{0,1\}$ & Classical Perceptron\\
\hline
Sigmoid & $(0,1)$ & Binary Classification Output Layer\\
\hline
Tanh & $(-1,1)$ & Hidden Layers (Traditional Networks)\\
\hline
ReLU & $[0,\infty)$ & Hidden Layers in CNNs and MLPs\\
\hline
Leaky ReLU & $(-\infty,\infty)$ & Avoiding Dying ReLU\\
\hline
Softmax & $(0,1)$ & Multi-class Classification Output Layer\\
\hline
\end{tabular}
\end{center}

\textbf{Note:} This experiment uses only the Step Activation Function. The remaining activation functions will be studied in subsequent experiments involving Multilayer Perceptrons and Deep Neural Networks.

%----------------------------------------------------------

\section*{5. Dataset Description}

\textbf{Dataset:} Banknote Authentication Dataset

\textbf{Source:} UCI Machine Learning Repository

\textbf{Download Link:}

\url{https://archive.ics.uci.edu/dataset/267/banknote+authentication}

\begin{center}
\begin{tabular}{|l|l|}
\hline
Instances & 1372\\
\hline
Features & 4 Numerical Features\\
\hline
Classes & 2\\
\hline
Missing Values & None\\
\hline
Task & Binary Classification\\
\hline
\end{tabular}
\end{center}

\textbf{Feature Description}

\begin{itemize}
\item Variance
\item Skewness
\item Curtosis
\item Entropy
\end{itemize}

Target Class

\begin{itemize}
\item 0 -- Authentic Banknote
\item 1 -- Forged Banknote
\end{itemize}

%----------------------------------------------------------

\section*{6. Experimental Procedure}

\textbf{Task 1: Dataset Exploration}

\begin{itemize}
\item Load the dataset.
\section{Load the dataset.}

\begin{lstlisting}
import numpy as np
import pandas as pd
import matplotlib.pyplot as plt
import seaborn as sns
from sklearn.model_selection import train_test_split
from sklearn.preprocessing import MinMaxScaler
from sklearn.metrics import accuracy_score, precision_score, recall_score
from sklearn.metrics import f1_score, confusion_matrix, ConfusionMatrixDisplay

\end{lstlisting}

\item Display the first five samples.
\section{Display the first five samples.}

\subsection{Python Code}

\begin{lstlisting}[language=Python]
columns = ["Variance","Skewness","Curtosis","Entropy","Class"]

df = pd.read_csv("data_banknote_authentication.txt",
                 names=columns)

print(df.head())
\end{lstlisting}

\subsection{Output}

\begin{verbatim}
   Variance  Skewness  Curtosis  Entropy  Class
0   3.62160    8.6661   -2.8073 -0.44699      0
1   4.54590    8.1674   -2.4586 -1.46210      0
2   3.86600   -2.6383    1.9242  0.10645      0
3   3.45660    9.5228   -4.0112 -3.59440      0
4   0.32924   -4.4552    4.5718 -0.98880      0
\end{verbatim}
\item Determine dataset dimensions.
\item Identify missing values.
\item Display descriptive statistics.
\section{Determine dataset dimensions,
Identify missing values,
Display descriptive statistics.}

\subsection{Python Code}

\begin{lstlisting}[language=Python]
print("Shape:",df.shape)

print("\nMissing Values")
print(df.isnull().sum())

print("\nDataset Info")
print(df.info())

print("\nStatistics")
print(df.describe())
\end{lstlisting}

\subsection{Output}

\begin{verbatim}
Shape: (1372, 5)

Missing Values
Variance    0
Skewness    0
Curtosis    0
Entropy     0
Class       0
dtype: int64

Dataset Info
<class 'pandas.core.frame.DataFrame'>
RangeIndex: 1372 entries, 0 to 1371
Data columns (total 5 columns):
 #   Column    Non-Null Count  Dtype  
---  ------    --------------  -----  
 0   Variance  1372 non-null   float64
 1   Skewness  1372 non-null   float64
 2   Curtosis  1372 non-null   float64
 3   Entropy   1372 non-null   float64
 4   Class     1372 non-null   int64  
dtypes: float64(4), int64(1)
memory usage: 53.7 KB
None

Statistics
          Variance     Skewness     Curtosis      Entropy        Class
count  1372.000000  1372.000000  1372.000000  1372.000000  1372.000000
mean      0.433735     1.922353     1.397627    -1.191657     0.444606
std       2.842763     5.869047     4.310030     2.101013     0.497103
min      -7.042100   -13.773100    -5.286100    -8.548200     0.000000
25%      -1.773000    -1.708200    -1.574975    -2.413450     0.000000
50%       0.496180     2.319650     0.616630    -0.586650     0.000000
75%       2.821475     6.814625     3.179250     0.394810     1.000000
max       6.824800    12.951600    17.927400     2.449500     1.000000

\end{verbatim}
\end{itemize}


\vspace{12mm}

\textbf{Task 2: Exploratory Data Analysis}

Generate
Histogram
\section{Histograms}

\subsection{Python Code}

\begin{lstlisting}[language=Python]
df.hist(figsize=(10,8))
plt.tight_layout()
plt.show()
\end{lstlisting}

Output

\begin{figure}[H]   % <---- CHANGED from [h] to [H]
    \centering
    \includegraphics[width=0.8\textwidth]{histogram}
    \caption{Histogram of the Dataset}
    \label{fig:histogram}
\end{figure}
\textit{Inference:} All features are skewed and spread out – no neat bell curves. Multiple peaks suggest the two classes are naturally separated.


Correlation Heatmap
\section{Correlation Heatmap}

\subsection{Python Code}
\begin{lstlisting}[language=Python]
plt.figure(figsize=(8,6))
sns.heatmap(df.corr(),
            annot=True,
            cmap="coolwarm")

plt.title("Correlation Heatmap")
plt.show()
\end{lstlisting}
Output
\begin{figure}[H]   % <---- CHANGED
    \centering
    \includegraphics[width=0.8\textwidth]{heatplot}
    \caption{Histogram of the Dataset}
    \label{fig:histogram}
\end{figure}

\textit{Inference:} No strong correlations between features – they're mostly independent. No redundant info to mess up the perceptron.
\vspace{3cm}
Scatter Plot
\section{Scatter Plot}

\subsection{Python Code}

\begin{lstlisting}[language=Python]
plt.figure(figsize=(7,6))

sns.scatterplot(
    x="Variance",
    y="Skewness",
    hue="Class",
    data=df
)

plt.show()
\end{lstlisting}

\subsection{Output}

\begin{figure}[H]   % <---- CHANGED
    \centering
    \includegraphics[width=0.8\textwidth]{scatterplot}
    \caption{Histogram of the Dataset}
    \label{fig:histogram}
\end{figure}
\textit{Inference:} The two classes form clear, separate clusters. They're roughly linearly separable – exactly why our perceptron works well.
\vspace{1cm}

Boxplots
\section{Box Plot}

\subsection{Python Code}

\begin{lstlisting}[language=Python]
plt.figure(figsize=(10,6))

for i,col in enumerate(df.columns[:-1]):
    plt.subplot(2,2,i+1)
    sns.boxplot(y=df[col])

plt.tight_layout()
plt.show()
\end{lstlisting}

\subsection{Output}

\begin{figure}[H]   % <---- CHANGED
    \centering
    \includegraphics[width=0.8\textwidth]{boxplot}
    \caption{Histogram of the Dataset}
    \label{fig:histogram}
\end{figure}
\textit{Inference:} Features have very different scales and quite a few outliers. Normalisation was a must – otherwise big numbers dominate weight updates.
\vspace{2mm}

\textbf{Task 3: Data Preprocessing}

\begin{itemize}
\item Normalize all numerical features.
\item Split the dataset into Training (80\%) and Testing (20\%).
\section{Normalizing all numerical features and Split Data}

\subsection{Python Code}

\begin{lstlisting}[language=Python]
X = df.iloc[:,:-1].values
y = df.iloc[:,-1].values

scaler = MinMaxScaler()

X = scaler.fit_transform(X)

X_train,X_test,y_train,y_test = train_test_split(
    X,
    y,
    test_size=0.2,
    random_state=42
)
\end{lstlisting}
\end{itemize}

\vspace{2mm}

\textbf{Task 4: Perceptron Implementation}

Implement from scratch

\begin{itemize}
\item Weight Initialization
\item Bias Initialization
\item Step Activation Function
\item Forward Propagation
\item Perceptron Learning Rule
\section{Perceptron}

\subsection{Python Code}

\begin{lstlisting}[language=Python]
import numpy as np

class Perceptron:

    def __init__(self, lr=0.01, epochs=20):
        self.lr = lr
        self.epochs = epochs
        self.weights = None
        self.bias = 0

    def activation(self, z):
        return np.where(z >= 0, 1, 0)

    def fit(self, X, y):

        n_features = X.shape[1]

        self.weights = np.zeros(n_features)
        self.bias = 0

        self.errors = []
        self.weight_history = []
        self.bias_history = []

        for epoch in range(self.epochs):

            errors = 0

            for xi, target in zip(X, y):

                z = np.dot(xi, self.weights) + self.bias
                prediction = self.activation(z)

                update = self.lr * (target - prediction)

                self.weights += update * xi
                self.bias += update

                if update != 0:
                    errors += 1

            self.errors.append(errors)
            self.weight_history.append(self.weights.copy())
            self.bias_history.append(self.bias)

            print(f"Epoch {epoch+1}: Errors = {errors}")

    def predict(self, X):
        z = np.dot(X, self.weights) + self.bias
        return self.activation(z)
\end{lstlisting}

\end{itemize}

\vspace{2mm}

\textbf{Task 5: Model Training}

Display

\begin{itemize}
\item Epoch Number
\item Misclassified Samples
\item Updated Weights
\item Updated Bias
\section{Train Model}

\subsection{Python Code}

\begin{lstlisting}[language=Python]
model = Perceptron(lr=0.01, epochs=20)
model.fit(X_train, y_train)

y_pred = model.predict(X_test)
\end{lstlisting}

\subsection{Output}

\begin{verbatim}
Epoch 1: Errors = 166
Epoch 2: Errors = 88
Epoch 3: Errors = 57
Epoch 4: Errors = 51
Epoch 5: Errors = 47
Epoch 6: Errors = 45
Epoch 7: Errors = 48
Epoch 8: Errors = 32
Epoch 9: Errors = 19
Epoch 10: Errors = 40
Epoch 11: Errors = 35
Epoch 12: Errors = 28
Epoch 13: Errors = 40
Epoch 14: Errors = 31
Epoch 15: Errors = 28
Epoch 16: Errors = 38
Epoch 17: Errors = 27
Epoch 18: Errors = 28
Epoch 19: Errors = 32
Epoch 20: Errors = 31

\end{verbatim}
\textit{Inference:} Errors dropped from 166 to around 30 in no time. Minor wobbles later – typical when data isn't perfectly separable.
\end{itemize}

\vspace{2mm}

\textbf{Task 6: Model Evaluation}

Compute

\begin{itemize}
\item Accuracy
\item Precision
\item Recall
\item F1-score
\section{Modal Evaluation}


\subsection{Python Code}
\begin{lstlisting}[language=Python]
y_pred=model.predict(X_test)

print("Accuracy :",accuracy_score(y_test,y_pred))

print("Precision :",precision_score(y_test,y_pred))

print("Recall :",recall_score(y_test,y_pred))

print("F1 Score :",f1_score(y_test,y_pred))
\end{lstlisting}


\subsection{Output}

\begin{verbatim}
Accuracy : 0.9745454545454545
Precision : 0.9477611940298507
Recall : 1.0
F1 Score : 0.9731800766283525

\end{verbatim}
\item Confusion Matrix
\section{Confusion Matrix}

\subsection{Python Code}
\begin{lstlisting}[language=Python]
cm=confusion_matrix(y_test,y_pred)

disp=ConfusionMatrixDisplay(cm)

disp.plot()

plt.show()
\end{lstlisting}


Output

\begin{figure}[H]   % <---- CHANGED
    \centering
    \includegraphics[width=0.8\textwidth]{confusion matrix.png}
    \caption{Histogram of the Dataset}
    \label{fig:Confusion Matrix}
\end{figure}
\textit{Inference:} Almost everything lies on the diagonal – very few misclassifications. The perceptron clearly nailed the real vs fake distinction.
\end{itemize}

%----------------------------------------------------------

\section*{7. Mandatory Plots}

\begin{longtable}{|p{3.4cm}|p{3.4cm}|p{6cm}|}
\hline
\textbf{Plot} & \textbf{Purpose} & \textbf{Expected Inference}\\
\hline
Feature Histograms & Distribution Analysis & Identify skewness and outliers.\\
\hline
Correlation Heatmap & Feature Relationship & Observe highly correlated features.\\
\hline
Scatter Plot & Class Distribution & Determine whether classes are approximately linearly separable.\\
\hline
Training Error vs Epoch & Learning Behaviour & Verify convergence of the perceptron.\\
\hline
Weight Evolution & Parameter Analysis & Observe stabilization of learned weights.\\
\hline
Bias Evolution & Parameter Analysis & Understand the role of the bias term.\\
\hline
Confusion Matrix & Performance Evaluation & Interpret TP, FP, TN and FN.\\
\hline
Learning Rate Comparison & Hyperparameter Study & Compare convergence behaviour for different learning rates.\\
\hline
Decision Boundary (Optional) & Visualization & Illustrate the separating hyperplane using two selected features.\\
\hline
\end{longtable}

Students should provide a brief interpretation (2--3 sentences) for every generated plot.

\section{Weight Evolution}

\subsection{Python Code}
\begin{lstlisting}[language=Python]
weights=np.array(model.weight_history)

for i in range(weights.shape[1]):
    plt.plot(weights[:,i],label=f"W{i+1}")

plt.legend()

plt.xlabel("Epoch")

plt.ylabel("Weight Value")

plt.title("Weight Evolution")

plt.grid()

plt.show()

\end{lstlisting}
Output
\begin{figure}[H]   
    \centering
    \includegraphics[width=0.8\textwidth]{Weight Evolution.png}
    \caption{Histogram of the Dataset}
    \label{fig:Weight Evolution}
\end{figure}
\textit{Inference:} Weights settle quickly after the first few epochs and stay steady. No wild oscillations – training was smooth and converged well.

Output


\section{Bias Evolution}

\subsection{Python Code}
\begin{lstlisting}[language=Python]
plt.plot(model.bias_history,
         marker='o')

plt.xlabel("Epoch")

plt.ylabel("Bias")

plt.title("Bias Evolution")

plt.grid()

plt.show()
\end{lstlisting}


\subsection{Output}

\begin{figure}[H]   % <---- CHANGED
    \centering
    \includegraphics[width=0.8\textwidth]{biasevolution.png}
    \caption{Histogram of the Dataset}
    \label{fig:Bias Evolution}
\end{figure}
\textit{Inference:} Bias jumps sharply at the start, then plateaus. It just shifts the decision boundary to the right place and stays put.

%----------------------------------------------------------

\section*{8. Performance Tables}

\textbf{Training Summary}

\begin{verbatim}
          Parameter                                Value
0      Dataset Size                                 1372
1  Train/Test Split                           1097 / 275
2     Learning Rate                                 0.01
3            Epochs                                   20
4     Final Weights  [-0.1464, -0.1322, -0.1521, 0.0161]
5        Final Bias                                 0.19
6          Accuracy                               0.9745
7         Precision                               0.9478
8            Recall                                  1.0
9          F1-score                               0.9732

\end{verbatim}

\vspace{3mm}

\textbf{Epoch-wise Learning}


%----------------------------------------------------------

\section*{9. Additional Tasks}

\begin{enumerate}[leftmargin=*]
\item Compare the Step and Sigmoid activation functions. Discuss why the Step Activation Function is not used in modern deep learning models.
\item Compare the implementation with Scikit-learn's Perceptron.
\item Repeat the experiment using learning rates 0.001, 0.01 and 0.1.

\section{Experiment Using Different Learning Rates}

\subsection{Python Code}
\begin{lstlisting}[language=Python]
rates=[0.001,0.01,0.1]

plt.figure(figsize=(8,6))

for lr in rates:

    p=Perceptron(
        lr=lr,
        epochs=20
    )

    p.fit(X_train,y_train)

    plt.plot(p.errors,label=f"LR={lr}")

plt.legend()

plt.xlabel("Epoch")

plt.ylabel("Errors")

plt.title("Learning Rate Comparison")

plt.grid()

plt.show()
\end{lstlisting}


\subsection{Output}
\begin{verbatim}
Epoch 1: Errors = 166
Epoch 2: Errors = 88
Epoch 3: Errors = 57
Epoch 4: Errors = 51
Epoch 5: Errors = 47
Epoch 6: Errors = 45
Epoch 7: Errors = 48
Epoch 8: Errors = 32
Epoch 9: Errors = 19
Epoch 10: Errors = 40
Epoch 11: Errors = 35
Epoch 12: Errors = 28
Epoch 13: Errors = 40
Epoch 14: Errors = 31
Epoch 15: Errors = 28
Epoch 16: Errors = 38
Epoch 17: Errors = 27
Epoch 18: Errors = 28
Epoch 19: Errors = 32
Epoch 20: Errors = 31
Epoch 1: Errors = 166
Epoch 2: Errors = 88
Epoch 3: Errors = 57
Epoch 4: Errors = 51
Epoch 5: Errors = 47
Epoch 6: Errors = 45
Epoch 7: Errors = 48
Epoch 8: Errors = 32
Epoch 9: Errors = 19
Epoch 10: Errors = 40
Epoch 11: Errors = 35
Epoch 12: Errors = 28
Epoch 13: Errors = 40
Epoch 14: Errors = 31
Epoch 15: Errors = 28
Epoch 16: Errors = 38
Epoch 17: Errors = 27
Epoch 18: Errors = 28
Epoch 19: Errors = 32
Epoch 20: Errors = 31
Epoch 1: Errors = 166
Epoch 2: Errors = 88
Epoch 3: Errors = 57
Epoch 4: Errors = 51
Epoch 5: Errors = 47
Epoch 6: Errors = 45
Epoch 7: Errors = 48
Epoch 8: Errors = 32
Epoch 9: Errors = 19
Epoch 10: Errors = 40
Epoch 11: Errors = 35
Epoch 12: Errors = 28
Epoch 13: Errors = 40
Epoch 14: Errors = 31
Epoch 15: Errors = 28
Epoch 16: Errors = 38
Epoch 17: Errors = 27
Epoch 18: Errors = 28
Epoch 19: Errors = 32
Epoch 20: Errors = 31


\end{verbatim}
\textit{Inference:} Errors dropped from 166 to around 30 in no time. Minor wobbles later – typical when data isn't perfectly separable.

\begin{figure}[H]   % <---- CHANGED
    \centering
    \includegraphics[width=0.8\textwidth]{Learning rate.png}
    \caption{Histogram of the Dataset}
    \label{fig:Learning rate comparision}
\end{figure}
\textit{Inference:} LR = 0.01 is the sweet spot – fast enough to learn, stable enough not to bounce. 0.1 is jumpy, 0.001 is painfully slow. Middle ground wins.
\item Explain why a Single Layer Perceptron cannot solve the XOR problem.
\item Study the effect of feature normalization on model convergence.
\end{enumerate}

%----------------------------------------------------------

\section*{10. Report Contents}

The report should include:

\begin{itemize}
\item Objective
\item Background Theory
\item Activation Functions
\item Dataset Description
\item Experimental Procedure
\item Source Code
\item Experimental Results
\item Generated Plots with Interpretation
\item Performance Tables
\item Discussion
\item Conclusion
\item References
\end{itemize}

%----------------------------------------------------------

\section*{11. References}

\begin{enumerate}
\item F. Rosenblatt, ``The Perceptron,'' \emph{Psychological Review}, 1958.
\item I. Goodfellow, Y. Bengio and A. Courville, \emph{Deep Learning}, MIT Press, 2016.
\item C. M. Bishop, \emph{Pattern Recognition and Machine Learning}, Springer, 2006.
\item S. Haykin, \emph{Neural Networks and Learning Machines}, Pearson, 2009.
\item UCI Machine Learning Repository -- Banknote Authentication Dataset.
\item Scikit-learn Documentation: Perceptron.
\end{enumerate}



\begin{center}
    {\Huge \textbf{Additional Tasks }}
\end{center}

\textit1.  Code the perceptron learning algorithm for OR, NOT, and AND gates using Python. Show the weights after each update. Plot the decision boundary after each weight update using Python built-in packages. (K4)

2. Code the perceptron learning algorithm for the XOR gate using Python. Show the weights after each update. Plot the decision boundary after each weight update using Python's built-in packages. Analyze and identify the pattern that prevents the algorithm from converging. (K4)


\section{Task-1}

\subsection{Python Code}
\begin{lstlisting}[language=Python]
class Perceptron:

    def __init__(self, lr=0.1, epochs=20):
        self.lr = lr
        self.epochs = epochs
        self.weights = None
        self.bias = 0

    def activation(self, x):
        return 1 if x >= 0 else 0

    def predict(self, X):
        linear = np.dot(X, self.weights) + self.bias
        return np.array([self.activation(i) for i in linear])

    def plot_boundary(self, X, y, title):

        # ---------- NOT Gate ----------
        if X.shape[1] == 1:

            plt.figure(figsize=(5,2))

            colors = ['red' if i==0 else 'blue' for i in y]

            plt.scatter(X[:,0], np.zeros(len(X)),
                        c=colors, s=120)

            if abs(self.weights[0]) > 1e-8:
                threshold = -self.bias/self.weights[0]
                plt.axvline(threshold,
                            color='black',
                            label='Decision Boundary')

            plt.yticks([])
            plt.xlim(-0.5,1.5)
            plt.grid(True)
            plt.title(title)
            plt.legend()
            plt.show()

            return

        # ---------- AND / OR / XOR ----------

        plt.figure(figsize=(5,5))

        colors = ['red' if i==0 else 'blue' for i in y]

        plt.scatter(X[:,0], X[:,1],
                    c=colors,
                    s=120)

        if abs(self.weights[1]) > 1e-8:

            x = np.linspace(-0.5,1.5,100)

            y_line = -(self.weights[0]*x+self.bias)/self.weights[1]

            plt.plot(x,y_line,
                     'k-',
                     label='Decision Boundary')

        plt.xlim(-0.5,1.5)
        plt.ylim(-0.5,1.5)
        plt.grid(True)
        plt.title(title)
        plt.legend()
        plt.show()

    def fit(self, X, y, gate_name):

        self.weights = np.zeros(X.shape[1])
        self.bias = 0

        update = 1

        print("\n===================================")
        print(f"Training Perceptron for {gate_name}")
        print("===================================")

        for epoch in range(self.epochs):

            error_count = 0

            for xi, target in zip(X, y):

                linear = np.dot(xi, self.weights) + self.bias

                pred = self.activation(linear)

                error = target - pred

                if error != 0:

                    self.weights += self.lr * error * xi
                    self.bias += self.lr * error

                    print(f"\nUpdate {update}")
                    print("Input :", xi)
                    print("Target :", target)
                    print("Prediction :", pred)
                    print("Weights :", self.weights)
                    print("Bias :", self.bias)

                    # Plot every update except limit XOR to first 10
                    if gate_name == "XOR":
                        if update <= 10:
                            self.plot_boundary(
                                X,
                                y,
                                f"{gate_name} Update {update}"
                            )
                    else:
                        self.plot_boundary(
                            X,
                            y,
                            f"{gate_name} Update {update}"
                        )

                    update += 1
                    error_count += 1

            if error_count == 0:
                print("\nConverged Successfully!")
                break

        print("\nFinal Weights :", self.weights)
        print("Final Bias :", self.bias)
        print("Predictions :", self.predict(X))
\end{lstlisting}
\subsection{AND}
\begin{lstlisting}[language=Python]
X = np.array([
    [0,0],
    [0,1],
    [1,0],
    [1,1]
])

y = np.array([0,0,0,1])

model = Perceptron(lr=0.1, epochs=20)
model.fit(X, y, "AND")
\end{lstlisting}

\begin{verbatim}
Update 1
Input : [0 0]
Target : 0
Prediction : 1
Weights : [0. 0.]
Bias : -0.1

Update 2
Input : [1 1]
Target : 1
Prediction : 0
Weights : [0.1 0.1]
Bias : 0.0

Update 3
Input : [0 0]
Target : 0
Prediction : 1
Weights : [0.1 0.1]
Bias : -0.1

Update 4
Input : [0 1]
Target : 0
Prediction : 1
Weights : [0.1 0. ]
Bias : -0.2

Update 5
Input : [1 1]
Target : 1
Prediction : 0
Weights : [0.2 0.1]
Bias : -0.1

Update 6
Input : [0 1]
Target : 0
Prediction : 1
Weights : [0.2 0. ]
Bias : -0.2

Update 7
Input : [1 0]
Target : 0
Prediction : 1
Weights : [0.1 0. ]
Bias : -0.30000000000000004

Update 8
Input : [1 1]
Target : 1
Prediction : 0
Weights : [0.2 0.1]
Bias : -0.20000000000000004

Converged Successfully!

Final Weights : [0.2 0.1]
Final Bias : -0.20000000000000004
Predictions : [0 0 0 1]
\end{verbatim}

\subsection{OR}
\begin{lstlisting}[language=Python]
X = np.array([
    [0,0],
    [0,1],
    [1,0],
    [1,1]
])

y = np.array([0,1,1,1])

model = Perceptron(lr=0.1, epochs=20)
model.fit(X, y, "OR")
\end{lstlisting}

\begin{verbatim}
Update 1
Input : [0 0]
Target : 0
Prediction : 1
Weights : [0. 0.]
Bias : -0.1


Update 2
Input : [0 1]
Target : 1
Prediction : 0
Weights : [0.  0.1]
Bias : 0.0


Update 3
Input : [0 0]
Target : 0
Prediction : 1
Weights : [0.  0.1]
Bias : -0.1


Update 4
Input : [1 0]
Target : 1
Prediction : 0
Weights : [0.1 0.1]
Bias : 0.0


Update 5
Input : [0 0]
Target : 0
Prediction : 1
Weights : [0.1 0.1]
Bias : -0.1


Converged Successfully!

Final Weights : [0.1 0.1]
Final Bias : -0.1
Predictions : [0 1 1 1]


\end{verbatim}

\subsection{NOR}
\begin{lstlisting}[language=Python]
X = np.array([
    [0],
    [1]
])

y = np.array([1,0])

model = Perceptron(lr=0.1, epochs=20)
model.fit(X, y, "NOT")
\end{lstlisting}

\begin{verbatim}
Update 1
Input : [1]
Target : 0
Prediction : 1
Weights : [-0.1]
Bias : -0.1


Update 2
Input : [0]
Target : 1
Prediction : 0
Weights : [-0.1]
Bias : 0.0


Converged Successfully!

Final Weights : [-0.1]
Final Bias : 0.0
Predictions : [1 0]

\end{verbatim}


\section{Task-2}

\subsection{XOR}
\begin{lstlisting}[language=Python]
X = np.array([
    [0,0],
    [0,1],
    [1,0],
    [1,1]
])

y = np.array([0,1,1,0])

model = Perceptron(lr=0.1, epochs=20)
model.fit(X, y, "XOR")
\end{lstlisting}

\begin{verbatim}
Update 1
Input : [0 0]
Target : 0
Prediction : 1
Weights : [0. 0.]
Bias : -0.1


Update 2
Input : [0 1]
Target : 1
Prediction : 0
Weights : [0.  0.1]
Bias : 0.0


Update 3
Input : [1 1]
Target : 0
Prediction : 1
Weights : [-0.1  0. ]
Bias : -0.1


Update 4
Input : [0 1]
Target : 1
Prediction : 0
Weights : [-0.1  0.1]
Bias : 0.0


Update 5
Input : [1 0]
Target : 1
Prediction : 0
Weights : [0.  0.1]
Bias : 0.1


Update 6
Input : [1 1]
Target : 0
Prediction : 1
Weights : [-0.1  0. ]
Bias : 0.0


Update 7
Input : [0 0]
Target : 0
Prediction : 1
Weights : [-0.1  0. ]
Bias : -0.1


Update 8
Input : [0 1]
Target : 1
Prediction : 0
Weights : [-0.1  0.1]
Bias : 0.0


Update 9
Input : [1 0]
Target : 1
Prediction : 0
Weights : [0.  0.1]
Bias : 0.1


Update 10
Input : [1 1]
Target : 0
Prediction : 1
Weights : [-0.1  0. ]
Bias : 0.0

an soooo.. On
\end{verbatim}

\subsection*{Analysis of XOR Non‑Convergence}

\begin{verbatim}
The weights keep oscillating between three states because every correction 
that fixes one misclassified point breaks another. This happens because 
XOR is not linearly separable - no straight line can separate (0,0) and 
(1,1) from (0,1) and (1,0). So the perceptron keeps chasing its tail and 
never converges.
\end{verbatim}



\end{document}